\pdfoutput=1
\documentclass[11pt]{article}
\usepackage[margin=1in]{geometry}
\usepackage{amsmath,amssymb,amsthm}
\usepackage{enumitem}
\usepackage{booktabs}
\usepackage{array}
\usepackage{hyperref}
\hypersetup{hidelinks,
  pdftitle={Scale preservation is not enough: a profile-level independence countermodel for the Corollary 3.12 assembly inputs},
  pdfauthor={Toshihisa Urahigashi}}

\newcommand{\uu}[1]{\underline{\underline{#1}}}
\newcommand{\Fam}{\mathrm{Fam}}
\newtheorem{theorem}{Theorem}[section]
\newtheorem{proposition}[theorem]{Proposition}
\newtheorem{corollary}[theorem]{Corollary}
\theoremstyle{definition}
\newtheorem{definition}[theorem]{Definition}
\newtheorem{remark}[theorem]{Remark}

\title{Scale preservation is not enough:\\
a profile-level independence countermodel\\
for the Corollary~3.12 assembly inputs}
\author{Toshihisa Urahigashi\\
\small\texttt{123026.urahigashi.toshihisa@gmail.com}}
\date{Version 1.0.0 --- 13 November 2026\\[3pt]
\small Independent researcher. This work was carried out personally and is
unrelated to any institutional affiliation.\\[3pt]
\small ORCID \href{https://orcid.org/0009-0004-0460-6242}{\texttt{0009-0004-0460-6242}}\\[-2pt]
\small \href{https://doi.org/10.5281/zenodo.24195284}{doi:10.5281/zenodo.24195284} (all versions: \href{https://doi.org/10.5281/zenodo.24195283}{10.5281/zenodo.24195283})}

\begin{document}
\maketitle

\begin{abstract}
We exhibit a concrete profile-valued countermodel showing that the
remaining value-group profile switch in a conditional assembly from
Theorem~3.11 to Corollary~3.12 in inter-universal Teichm\"uller theory
is logically independent of the other assembly inputs.  At
$\ell=7$, with labels $j\in\{1,2,3\}$ at one bad place and
normalized unit $u=1$, the theta-side readouts are $-j^2$ and the
$q$-side readouts are $-1$, with equal procession weights $1/3$.
The normalized generator and the unit-side readouts agree;
membership, monotonicity, the hull domain and both volume
decompositions hold.  Nevertheless the aggregates are $-14/3$ and
$-1$, the aggregate scale is forced to $3/14$, and the
Corollary-3.12-shaped conclusion fails.  For general odd primes
$\ell\ge5$ the aggregate scale is $12/(\ell(\ell+1))$.  The
countermodel and the resulting independence corollary are
formalized in Lean~4; in the recorded selected build they depend
only on the standard three axioms.

The profile and normalizations are not arbitrary choices.  They are
the source readouts: theta values $\uu q^{\,j^2}$, plain procession
averaging, and the coefficient comparison in the proof of
[IUTchIV, Theorem~1.10].  This note does not claim that
Mochizuki's actual construction is this model, nor that
Theorem~3.11, Corollary~3.12, the ABC conjecture or the switch
H*val is false.  It proves only that no derivation exists from the
listed inputs once labelwise profile equality is removed.  Further
progress therefore requires a genuinely profile-level input, or a
richer formalization of the actual construction, rather than more
lemmas inside the same abstraction.

\medskip
\noindent\textbf{Keywords:} inter-universal Teichm\"uller theory,
Corollary~3.12, value-group profile, scale preservation,
independence countermodel, Lean.
\end{abstract}

\section*{What is and is not claimed}

\textbf{Claimed.}
\begin{itemize}[leftmargin=*]
\item The hypotheses of the assembly capstone, with the single
labelwise equality hypothesis removed, do not imply the
Corollary-3.12-shaped volume conclusion.  A concrete datum
satisfies all remaining hypotheses and falsifies the conclusion.
\item In this countermodel the chart/generator scale and the unit
side are preserved; only the labelwise profile differs.  The
aggregate scale is nevertheless forced away from $1$.
\item The numerical profile, equal weights and $q$-side
normalization used below are read off from the definitions and
computations in IUT~III--IV, as documented in Section~3.
\end{itemize}

\textbf{Not claimed.}
\begin{itemize}[leftmargin=*]
\item That Mochizuki's actual Frobenioid construction reduces to
the simple model in Section~4.
\item That Theorem~3.11, Corollary~3.12, the ABC conjecture or
H*val is false.
\item That the whole assembly paper, the IUT black boxes, or frame
fidelity have been formalized in Lean.
\end{itemize}

\section{Scale versus profile}

\subsection{Two meanings of scale preservation}

The phrase ``the scale is preserved'' has been used at two
different layers in the assembly from Theorem~3.11 to
Corollary~3.12.  The first layer concerns calibration of the
measuring device itself: a chart identifying two presentations has
distortion $s_{\mathrm{chart}}=1$, or a normalized generator has
the same readout on both sides.  The second layer concerns the
aggregate obtained after summing labelwise components of the
measured object with procession weights.

Fix one bad place and a positive normalized unit $u$.  Write
$\theta_j$ and $q_j$ for the label-$j$ theta-side and $q$-side
readouts, and $w_j$ for the weights.  The aggregate ratio is
\[
s_{\mathrm{agg}}
  :=\frac{\sum_j w_jq_j}{\sum_jw_j\theta_j}.
\]
Even when the chart and the weights are the same on both sides,
$s_{\mathrm{agg}}$ need not equal $1$ if the sequences
$\{\theta_j\}_j$ and $\{q_j\}_j$ differ.  We call such a sequence
a \emph{profile}.  A profile is vector-valued information which
becomes a scalar only after aggregation.

\subsection{Relation to previous work}

The note [U26] showed, by a countermodel, that
Corollary~3.12 does not follow from the statements actually cited
in its proof alone.  The subsequent assembly note [U26c] used
material outside that cited range to supply the scalar and unit
layers at the statement level, while isolating the value-group
profile switch H*val (there called R2V) as the central remaining
input.

The present note connects these two results.  The aggregate scale
$3/14$ in the countermodel of [U26] was described at the level of
that cited theory as a scalar defect.  At the level of concrete
values it is forced by averaging the $j^2$ profile.  Here the same
phenomenon is placed inside the hypothesis shape of the assembly
capstone of [U26c].  All supplied layers are satisfied, the
profile alone differs, and the conclusion still fails.  Thus the
result is not a repetition of [U26]: it localizes the missing
input inside the expanded assembly.

\subsection{The question}

We answer the following question in the negative: assuming a
monotone volume, the hull domain condition, membership in the
family of possible images, finite-sum decompositions on both
sides, agreement at the normalized generator, and agreement of
unit-side readouts, does the Corollary-3.12-shaped volume
conclusion follow without labelwise profile equality?

\section{The abstract assembly shape}

We use the following abstraction of the capstone
\texttt{cor312\_of\_assembly} in the Lean assembly development.
Let $\Omega$ be a hull operator on a type $\alpha$, let
$\Fam\subseteq\mathcal P(\alpha)$ be a family of possible images,
let $T\subseteq\alpha$, let
$\mu:\mathcal P(\alpha)\to\mathbb R$ be a volume readout, and let
$v_Q\in\mathbb R$.  Let $I$ be a finite index set with weights
$w_i$ and readouts $\theta_i,q_i$.  Consider the inputs
\begin{center}
\begin{tabular}{lll}
\toprule
 & input & in the model below\\
\midrule
(A1) & $\mu$ is monotone & holds\\
(A2) & $\bigcup\Fam$ lies in the domain of $\Omega$ & holds\\
(A3) & $T\in\Fam$ & holds\\
(A4) & $\mu(T)=\sum_{i\in I}w_i\theta_i$ & both sides $-14/3$\\
(A5) & $v_Q=\sum_{i\in I}w_iq_i$ & both sides $-1$\\
(G)  & normalized generator readouts agree & both $-1$ at $j=1$\\
(U)  & unit-side readouts agree labelwise & both $0$\\
\bottomrule
\end{tabular}
\end{center}
The hypothesis deliberately removed is
\[
\tag{Hloc}
\forall i\in I,\qquad q_i=\theta_i .
\]
The conclusion under test is the Corollary-3.12-shaped comparison
\[
\tag{C}
v_Q\leq
\mu\!\left(\Omega.\varphi\left(\bigcup\Fam\right)\right).
\]
In the usual assembly one obtains $\mu(T)\leq\mu(\mathrm{hull})$
from membership and monotonicity (using extensivity of the hull
operator), and $v_Q=\mu(T)$ from (Hloc), (A4) and (A5).  The
countermodel shows that (Hloc) is genuinely load-bearing in this
inference.

The finite set $I$ is arbitrary at this abstract level.  Below we
fix one bad place and take $I$ to be its label set.  Reading $I$
instead as a set of places pre-aggregated over labels does not
remove the mismatch: at the fixed bad place the two aggregates
are already $-14/3$ and $-1$ (Section~7.2).

\section{The source readout}

An independence countermodel can be built from arbitrary
sequences.  What matters here is that the sequence, weights and
normalization below are those read off from the source.  This
section lists only the inputs needed to fix the numbers; it is
not a reproof of the source construction.

\subsection{Labels and the theta pilot}

For an odd prime $\ell$, the label set is
$\mathbb F_\ell^{\divideontimes}
 =\mathbb F_\ell^\times/\{\pm1\}$, with representatives
$j\in\{1,\ldots,\ell^\divideontimes\}$,
$\ell^\divideontimes=(\ell-1)/2$.  At a bad place the components
of the theta pilot occur as the system of normalized theta values
$\{\uu q_{\underline v}^{\,j^2}\}_j$ ([IUTchIII], Introduction,
Definition~3.8(i), Remark~3.11.1(i)).  Here $\uu q$ is not the
bare Tate parameter but the normalized generator of the next
paragraph.

\subsection{The $q$-side unit}

[IUTchIV, Theorem~1.10] puts
\[
\left|\log(\uu q)\right|
  =\frac1{2\ell}\log(\mathfrak q).
\]
We write $u_{\underline v}:=|\log(\uu q_{\underline v})|>0$ and,
with the sign convention for log-volumes, label readouts
$-j^2u_{\underline v}$ on the theta side and
$-u_{\underline v}$ on the $q$ side.  In the main countermodel
one bad place is fixed and $u_{\underline v}$ is normalized to
$1$.

\subsection{Capsules and cancellation of the tensor factor}

In the procession normalization in the proof of
[IUTchIV, Theorem~1.10], one evaluates the components of the
capsule $S^\pm_{j+1}=\{0,1,\ldots,j\}$ and then averages over
$j\in\{1,\ldots,\ell^\divideontimes\}$.  The distinguished factor
is $i^\dagger=j$, and at a bad place the degree entering there is
$\lambda=\operatorname{ord}(\uu q_{\underline v_j}^{\,j^2})$.
A factor $j+1$ occurs in the per-factor quantities, but after the
weighted average the capsule main term is
$-\frac{j^2}{2\ell}\log(\mathfrak q)$: the number of tensor
factors cancels.  The procession-averaged main term is
$-\frac{\ell+1}{24}\log(\mathfrak q)$.  Comparing coefficients
relative to $\log(\mathfrak q)$ gives
\[
\text{theta coefficient}=\frac{\ell+1}{24},
\qquad
\text{$q$ coefficient}=\frac1{2\ell},
\qquad
\frac{\text{theta}}{\text{$q$}}=\frac{\ell(\ell+1)}{12}.
\]
The reciprocal aggregate scale is $12/(\ell(\ell+1))$; at
$\ell=7$ this is $3/14$, as in [U26].

\subsection{Equal weights and the determinant}

Procession normalization in [IUTchIII, Proposition~3.9(i)] is
defined as the plain average over capsules.  Remark~3.9.3
explicitly answers ``no'' to the question whether more general
weighted averages may be used: compatibility with capsule
inclusions forces the same weight for each capsule.  Moreover
$\det_{\otimes M}$ is normalized so as to reproduce the packet
log-volume and may be omitted in volume-only arguments; the
$M$-th tensor power comparison of (Ob4) is a common factor on
both sides and does not change the profile ratio.  The equal
weights $1/3$ at $\ell=7$ below are therefore the source
definition, not a simplifying choice.

\begin{center}
\begin{tabular}{lll}
\toprule
item & source location & value in the model\\
\midrule
labels &
$\ell^\divideontimes=(\ell-1)/2$ &
$\ell=7$, $\{1,2,3\}$\\
theta components &
$\{\uu q^{\,j^2}\}_j$ &
$-1,-4,-9$\\
$q$ components &
normalized triangle generator &
$-1,-1,-1$\\
weights &
Prop.~3.9(i), Rem.~3.9.3 &
$1/3,1/3,1/3$\\
aggregates &
Thm.~1.10 main terms &
$-14/3$ versus $-1$\\
\bottomrule
\end{tabular}
\end{center}

\section{The countermodel}

We transcribe the definitions of the Lean file listed in
Section~6.

\begin{definition}[Profile datum at $\ell=7$]
\begin{itemize}[leftmargin=*]
\item Index set: $\texttt{placesProf}=\{1,2,3\}$.
\item Weights: $w_j=1/3$ (\texttt{wProf}).
\item Theta readouts: $\theta_j=-j^2$ (\texttt{thetaVProf}).
\item $q$ readouts: $q_j=-1$ (\texttt{qVProf}).
\item Unit readouts: $0$ on both sides for every label
(\texttt{unitThetaProf}, \texttt{unitQProf}).
\item Carrier type $\alpha=\mathbb R$; hull operator $\Omega$ is
the identity hull \texttt{trivialHull}: every subset is in the
domain and is a hull, and $\Omega.\varphi(S)=S$.
\item Theta region and family:
$T=\{-14/3\}$ (\texttt{TProf}), $\Fam=\{T\}$ (\texttt{FamProf}).
\item Aggregate $q$ value: $v_Q=-1$ (\texttt{vQProf}).
\item Volume readout, for $S\subseteq\mathbb R$:
\[
\mu(S)=
\begin{cases}
-\dfrac{14}{3}, & -\dfrac{14}{3}\in S,\\[2mm]
-100, & -\dfrac{14}{3}\notin S
\end{cases}
\qquad\text{(\texttt{muProf})}.
\]
\end{itemize}
\end{definition}

The definitions of $T$ and $\mu$ are intentionally minimal.
They are not a reconstruction of the actual log-volume or
holomorphic hull.  Their purpose is to test what follows from
the weak inputs listed in Section~2 without additional geometry
of the actual hull.  If the actual construction has stronger
hull properties which rescue the conclusion, those properties
are precisely an extra input beyond the list; the present
independence result identifies that such an extra input is
necessary.

\section{Main theorems}

\begin{theorem}[Profile countermodel;
Lean: \texttt{profile\_countermodel}]
The datum of Definition~4.1 satisfies all of the following.
\begin{enumerate}[leftmargin=*]
\item $\mu$ is monotone, $\bigcup\Fam$ is in the domain of the
identity hull, and $T\in\Fam$.
\item The decompositions
\[
\mu(T)=\sum_{j=1}^{3}\frac13(-j^2)=-\frac{14}{3},
\qquad
v_Q=\sum_{j=1}^{3}\frac13(-1)=-1
\]
hold.
\item At the normalized generator, $\theta_1=q_1=-1$, and the
unit-side readouts agree at every label.
\item $v_Q=\frac{3}{14}\mu(T)$.
\item Nevertheless (Hloc) is false, and the conclusion
$v_Q\leq\mu(\Omega.\varphi(\bigcup\Fam))$ is false.
\end{enumerate}
\end{theorem}

\begin{proof}
For monotonicity, let $A\subseteq B$.  If $-14/3\in A$ then
$-14/3\in B$ and both values are $-14/3$.  If
$-14/3\notin A$, then $\mu(A)=-100$, while $\mu(B)$ is either
$-14/3$ or $-100$; hence $\mu(A)\leq\mu(B)$.  For the identity
hull every set is in the domain and $\Omega.\varphi(S)=S$.
Since $\Fam=\{T\}$, we have $T\in\Fam$ and
$\bigcup\Fam=T$.

The finite sums are direct:
\[
\sum_{j=1}^{3}\frac13(-j^2)
 =-\frac{1+4+9}{3}=-\frac{14}{3},
\qquad
\sum_{j=1}^{3}\frac13(-1)=-1.
\]
Also $-14/3\in T$, so $\mu(T)=-14/3$.  At the generator,
$\theta_1=-1^2=-1=q_1$, and both unit readouts are $0$ by
definition.  The aggregate relation is
$\frac{3}{14}\cdot(-14/3)=-1=v_Q$.

On the other hand, at $j=2$ we have $q_2=-1$ and
$\theta_2=-4$, so (Hloc) fails.  The conclusion would read
$-1=v_Q\leq\mu(\Omega.\varphi(T))=\mu(T)=-14/3$, which is
false.
\end{proof}

\begin{corollary}[Aggregate scale;
Lean: \texttt{aggregate\_scale\_eq},
\texttt{aggregate\_scale\_ne\_one}]
In the datum of Definition~4.1, any real number $s$ satisfying
$v_Q=s\,\mu(T)$ is equal to $3/14$.  In particular $s\neq1$.
\end{corollary}

\begin{proof}
Solve $-1=s(-14/3)$.
\end{proof}

\begin{theorem}[Independence;
Lean: \texttt{no\_cor312\_derivation\_without\_hlocal}]
There is no universally quantified derivation of the conclusion
(C) from inputs (A1)--(A5) alone.  Precisely, it is false that
for every hull operator $\Omega$, family $\Fam$, region $T$,
readout $\mu$, value $v_Q$, finite index set and weights and
readouts, inputs (A1)--(A5) imply (C).
\end{theorem}

\begin{proof}
Apply such a putative derivation to Definition~4.1.  By
Theorem~5.1 all hypotheses (A1)--(A5) hold while (C) is false,
a contradiction.
\end{proof}

\subsection{The computation for general $\ell$}

The formalized witness is at $\ell=7$; one counterexample refutes
the universal derivation.  The aggregate computation is the same
for every odd prime $\ell\ge5$.  Put $m=(\ell-1)/2$, labels
$j=1,\ldots,m$, equal weights $1/m$, $\theta_j=-j^2u$ and
$q_j=-u$ with $u>0$.  Then
\[
\frac1m\sum_{j=1}^{m}j^2
 =\frac{(m+1)(2m+1)}6
 =\frac{\ell(\ell+1)}{12},
\]
hence
\[
s_{\mathrm{agg}}
 =\frac{-u}{-\frac{\ell(\ell+1)}{12}u}
 =\frac{12}{\ell(\ell+1)}\neq1
\qquad(\ell\ge5).
\]
For example, at $\ell=5,7,11,13$ the mean profile coefficients
are $5/2$, $14/3$, $11$ and $91/6$.  The general computation is
the elementary finite-sum formula; the Lean theorem mechanizes
the $\ell=7$ witness.

\subsection{Permutation-invariant readouts}

Extracting only the label $j=1$ would give the generator value
$-u$.  But a procession volume readout does not distinguish
labels in this way.  In general, a permutation-invariant linear
readout must have equal weights: invariance under all
transpositions forces all weights to coincide, and normalization
on the constant profile $(1,\ldots,1)$ forces each weight to be
$1/m$.  Such a readout is therefore the mean, which returns
$\ell(\ell+1)/12\neq1$ on the theta side.  The minimum readout
does return $1$, but it is not additive over tensor factors and
hence is not a log-volume readout: for instance
$\min(1,4)+\min(4,1)=2$, whereas
$\min((1,4)+(4,1))=5$.

\section{Lean verification}

The Lean module and namespace are
\begin{verbatim}
M3.relative.ActualValuePartProfileIndependence
Kouchiku.M3.ActualValuePartProfileIndependence
\end{verbatim}
Its
main declarations are
{\small\texttt{profile\_countermodel}},
{\small\texttt{aggregate\_scale\_ne\_one}} and
{\small\texttt{no\_cor312\_derivation\_\allowbreak
without\_hlocal}}.
The selected build in the author's Lean development repository
was run as follows.
\begin{verbatim}
python3 tools/lean_selected.py \
  M3.relative.ActualValuePartProfileIndependence
\end{verbatim}
The terminal output observed was:
\begin{verbatim}
Selected closure: 2 modules
[1/2] M3.theorem311.Remark395HullIsAClosure [cached]
[2/2] M3.relative.ActualValuePartProfileIndependence
'...profile_countermodel' depends on axioms:
  [propext, Classical.choice, Quot.sound]
'...aggregate_scale_ne_one' depends on axioms:
  [propext, Classical.choice, Quot.sound]
'...no_cor312_derivation_without_hlocal' depends on axioms:
  [propext, Classical.choice, Quot.sound]
PASS: 2 modules.
\end{verbatim}
Thus there is no \texttt{sorryAx} and no additional axiom in the
main theorems.  The Lean file is recorded in commit
\texttt{bc25c4c9} of the author's development branch
\texttt{muse/valuepart-radial-20261110}, and is included in the
supplementary package for this note.

\begin{remark}[Scope of the build]
This selected build does not recompile the whole assembly
development.  The present file deliberately imports only the
defining module of the \texttt{HullOperator} type also used by
the assembly capstone, together with the needed Mathlib
components.  The hypothesis shape was matched by direct reading
of \texttt{cor312\_of\_assembly} in
\texttt{Cor312Assembly.lean}.  What is machine-checked here is
therefore the independence theorem of this note itself, not the
assembly paper as a whole and not the actual IUT construction.
\end{remark}

\section{Objections and scope}

\subsection{The model is too degenerate}

The model uses the identity hull and a singleton region.  For an
independence statement quantified over all data satisfying the
listed properties, degenerate data are legitimate
counterexamples.  The actual holomorphic hull may of course have
stronger closure, convexity or analytic properties which rescue
the conclusion.  In that case the rescuing property is exactly
an additional input.  The present result certifies that no
derivation exists without such an addition.

\subsection{The indices should be places, not labels}

The finite sum in the capstone shape can be read with place
indices or with combined (place, label) indices.  We use the
finer reading to expose the profile.  Pre-aggregation at a fixed
bad place does not help: the two place aggregates are already
$-14u/3$ and $-u$.  Across several bad places the same ratio
$\ell(\ell+1)/12$ occurs at each place, and summing with
coefficients $u_{\underline v}\ge0$ produces no cancellation.
Coarsening the indices or passing to a global sum therefore does
not remove the mismatch.

\subsection{The readout $\mu$ is not the actual log-volume}

Correct, and stated as such above.  Inputs (A1)--(A5) impose on
$\mu$ only monotonicity and two decompositions.  Additional laws
satisfied by the actual log-volume --- its full behaviour under
tensor products, lower bounds coming from the spread of the
hull, or detailed normalizations across places --- are not in
the input list.  An argument using such laws may evade the
countermodel, but then the additional law carrying the weight
must be named.  The point of this note is that one cannot say
``the remaining lemmas will close the gap'' without that naming.

\subsection{Why not select the generator $j=1$?}

The two sides do agree at the normalized generator $j=1$.  That
is chart-scale preservation.  But the Corollary-3.12-shaped
conclusion measures the hull of the whole family of possible
images with procession-normalized volume, not a single
generator.  By Section~5.2, a permutation-invariant linear
aggregation is necessarily the mean.  An operation selecting
one label asymmetrically is not part of the mono-analytic
procession vocabulary; if additional structure justifies such a
selection, it forms part of the profile switch and should be
stated as such.

\subsection{Why not assume aggregate scale $s=1$ separately?}

Assuming $s_{\mathrm{agg}}=1$ removes the defect by fiat, but
under the source readout the two aggregates have ratio
$\ell(\ell+1)/12$, so $s_{\mathrm{agg}}=1$ is incompatible with
the readout itself.  The mathematical question is how the
theta-side quantity, before aggregation, is connected through
the hull of possible images to the $q$-side quantity.  Setting
a numerical ratio to $1$ afterwards does not prove that
connection.

\subsection{The actual construction may contain more}

This is compatible with the present conclusion.  For example,
[EtTh] contains a local skeleton in which the difference of two
sections is valued in the cyclotomic kernel and the theta class
is pinned modulo units.  The actual Frobenioid construction has
far more structure than the abstraction of Section~2.  What the
countermodel shows is that at least one piece of that structure
must carry weight at the profile level.  Conversely, once a
candidate additional mechanism is proposed, the countermodel
indicates which listed hypothesis it must violate or strengthen:
for instance the actual hull must spread strictly more than the
identity hull, the actual family of possible images must be
larger, or an object-level identification must change the
labelwise readouts.

\section{Conclusion}

Preservation of chart calibration and of the normalized
generator scale does not imply preservation after aggregation.
When the measured profile changes from $\{j^2u\}_j$ to
$\{u\}_j$, the same ruler and the same equal weights produce
aggregates differing by the factor $\ell(\ell+1)/12$.  The
countermodel satisfies every other assembly input and still
falsifies the conclusion.  The missing input is therefore
localized exactly at the profile switch, not distributed over
unspecified ``remaining lemmas''.

From the point of view of formalization, this is a boundary
result rather than a failure of effort.  A method which stacks
further lemmas inside the fixed abstraction cannot cross the
boundary.  Progress requires either formalizing the actual
construction in richer definitions and exhibiting where the
profile-level identification occurs, or explicitly introducing
an object-level input equivalent to that identification.  This
note fixes that necessity in finite, machine-checkable form.

\begin{center}
\begin{tabular}{lll}
\toprule
residual & content & relation to this note\\
\midrule
R1 & IUT black boxes & not verified here\\
R2V & value-group switch H*val &
central input shown independent\\
R2 & inhabitation / frame fidelity &
no inhabitation claim is made\\
R3 & independent re-verification &
source audit and rereading remain\\
\bottomrule
\end{tabular}
\end{center}

\section*{Disclosure}

This work was carried out by the author as an independent
researcher and is unrelated to any institutional affiliation.
Generative AI was used under the author's direction in the
audit, formalization and preparation of this draft.  The author
bears responsibility for its contents.  This is an unrefereed
preprint.

\begin{thebibliography}{99}

\bibitem[IUTchI]{IUTchI}
S.~Mochizuki,
Inter-universal Teichm\"uller Theory I: Construction of Hodge
Theaters, Publ.\ Res.\ Inst.\ Math.\ Sci.\ \textbf{57} (2021),
3--207.

\bibitem[IUTchII]{IUTchII}
S.~Mochizuki,
Inter-universal Teichm\"uller Theory II:
Hodge-Arakelov-theoretic Evaluation,
Publ.\ Res.\ Inst.\ Math.\ Sci.\ \textbf{57} (2021), 209--401.

\bibitem[IUTchIII]{IUTchIII}
S.~Mochizuki,
Inter-universal Teichm\"uller Theory III: Canonical Splittings
of the Log-theta-lattice,
Publ.\ Res.\ Inst.\ Math.\ Sci.\ \textbf{57} (2021), 403--626.

\bibitem[IUTchIV]{IUTchIV}
S.~Mochizuki,
Inter-universal Teichm\"uller Theory IV: Log-volume
Computations and Set-theoretic Foundations,
Publ.\ Res.\ Inst.\ Math.\ Sci.\ \textbf{57} (2021), 627--723.

\bibitem[EtTh]{EtTh}
S.~Mochizuki,
The \'Etale Theta Function and its Frobenioid-theoretic
Manifestations,
Publ.\ Res.\ Inst.\ Math.\ Sci.\ \textbf{45} (2009), 227--349.

\bibitem[SS]{SS}
P.~Scholze and J.~Stix, Why abc is still a conjecture, 2018.

\bibitem[U26]{U26}
T.~Urahigashi, Where Theorem 3.11 sits, Zenodo,
doi:10.5281/zenodo.24104932 (version 1.2.1).

\bibitem[U26b]{U26b}
T.~Urahigashi, Forgetting and Completeness, Zenodo,
doi:10.5281/zenodo.24165837 (version 1.1.3).

\bibitem[U26c]{U26c}
T.~Urahigashi, A conditional assembly from Theorem~3.11 to
Corollary~3.12: decomposition of the scalar layer, a semantic
principle, and isolation of the value-group profile switch,
Zenodo, doi:10.5281/zenodo.24179190 (version 1.0.1).

\bibitem[LANA]{LANA}
Project LANA, interim reports ((9-1), \texttt{xi\_iff}).
Cited only for the correspondence with preceding reductions;
no result of that project is used as an input here.

\end{thebibliography}

\end{document}




